% 集合与常用记号
\newcommand{\R}{\mathbb{R}}
\newcommand{\C}{\mathbb{C}}
\newcommand{\Z}{\mathbb{Z}}
\newcommand{\N}{\mathbb{N}}
\newcommand{\iu}{\mathrm{j}}
\newcommand{\e}{\mathrm{e}}
\newcommand{\dd}{\mathop{}\!\mathrm{d}}

% 信号与系统统一记号
\newcommand{\impulse}{\delta}
\newcommand{\stepfun}{u}
\newcommand{\conv}{\mathbin{\ast}}
\newcommand{\Sa}{\operatorname{Sa}}
\newcommand{\sinc}{\operatorname{sinc}}

% 变换算子；具体正、逆变换约定应在首次使用处声明。
\newcommand{\FT}{\mathcal{F}}
\newcommand{\IFT}{\mathcal{F}^{-1}}
\newcommand{\LT}{\mathcal{L}}
\newcommand{\ILT}{\mathcal{L}^{-1}}
\newcommand{\ZT}{\mathcal{Z}}
\newcommand{\IZT}{\mathcal{Z}^{-1}}
\newcommand{\HT}{\mathcal{H}}
\newcommand{\pv}{\operatorname{p.v.}}

% ── 重点标注 ────────────────────────────────────────────────
% 四类标注的分工、密度与写法约定见 docs/style-guide.md「重点标注」。
% 共同约定：只标文字短语，不标公式，标记内不出现 $...$ 或数学命令；
% 标记之间不嵌套、不重叠，同一短语只归一类。

% 红：结论与判据。读完后应当记住的知识要点。
\newcommand{\red}[1]{\textcolor{red}{#1}}

% 蓝：前提与检查项。结论成立的条件、适用范围，以及使用前应做的核对。
\definecolor{annotblue}{RGB}{10,80,170}
\newcommand{\blue}[1]{\textcolor{annotblue}{#1}}

% 加粗：术语与名称。正式名词、被定义对象、变换与性质的名称，只在首次出现处。
\newcommand{\keyterm}[1]{\textbf{#1}}

% 黄底：做题步骤与口诀。手上照做的计算顺序与速判口令。
% 逐字成盒，字间留出断行点，色块可随正文自然断行，不会顶出行宽。
\ExplSyntaxOn
\newcommand{\hl}[1]{%
  \penalty0
  \group_begin:
    \setlength{\fboxsep}{0pt}
    \colorbox{yellow!55}{\vphantom{国}\hspace*{1.2pt}}
    \penalty0
    \tl_map_inline:nn {#1} {
      \colorbox{yellow!55}{\vphantom{国}##1}
      \penalty0
    }
    \colorbox{yellow!55}{\vphantom{国}\hspace*{1.2pt}}
  \group_end:
  \penalty0}
\ExplSyntaxOff

% 成对关系，便于统一全文排版。
\newcommand{\fourierpair}{\mathrel{\xleftrightarrow{\,\FT\,}}}
\newcommand{\laplacepair}{\mathrel{\xleftrightarrow{\,\LT\,}}}
\newcommand{\zpair}{\mathrel{\xleftrightarrow{\,\ZT\,}}}
